% Unit 3 map -- the heaviest unit. Seven separate Zumdahl locations.
\documentclass{apchem-worksheet}

\apcunit{3}
\apcunittitle{Properties of Substances and Mixtures}
\apcdoctitle{Unit Map}
\apcsubtitle{CED Unit 3 \textbullet\ \textbf{18--22\% of the AP Exam} \textbullet\ 8 blocks + exam}

\begin{document}
\apctitleblock

\begin{apcnote}[Why this unit gets more time than any other]
Unit 3 is the largest slice of the AP exam --- roughly \textbf{one point in
five}. It is also the most scattered: its 13 topics live in seven different
places in Zumdahl, including an appendix. Follow the page ranges in the
table below rather than reading a chapter straight through; ``read chapter
10'' would give you barely a third of this unit.
\end{apcnote}

\section*{\sffamily Block calendar}

\renewcommand{\arraystretch}{1.3}
\noindent\begin{tabular}{@{}p{2.1cm}p{4.9cm}p{1.5cm}p{3.9cm}p{1.9cm}@{}}
\toprule
\textbf{Date} & \textbf{Topics} & \textbf{CED} & \textbf{Read (PDF pp.)} & \textbf{Practice}\\
\midrule
Fri Sep 11 & Intermolecular forces; polarizability & 3.1 & \S10.1 \textbullet\ 493--498 & WS 1\\
Mon Sep 14 & Solid types; phase changes; vapor pressure & 3.2, 3.3 & \S10.2--10.7 \textbullet\ 498--523 & WS 1\\
Wed Sep 16 & Gas laws; partial pressures & 3.4 & \S5.2--5.5 \textbullet\ 229--251 & WS 2\\
Fri Sep 18 & Kinetic molecular theory & 3.5 & \S5.6--5.7 \textbullet\ 251--260 & WS 2\\
Mon Sep 21 & Real gases; deviation from ideality & 3.6 & \S5.8--5.9 \textbullet\ 260--264 & WS 2\\
Wed Sep 23 & Solutions; molarity; particulate representations & 3.7, 3.8 & \S4.1--4.3 \textbullet\ 169--182; \S11.1--11.3 \textbullet\ 545--563 & WS 3\\
Fri Sep 25 & Solubility; separation techniques & 3.9, 3.10 & \S1.10 \textbullet\ 53--57; \S11.2 \textbullet\ 551--559 & WS 4\\
Mon Sep 28 & Spectroscopy; photons; Beer--Lambert & 3.11--3.13 & \S7.1--7.2 \textbullet\ 332--340; \textbf{App.~3} \textbullet\ 1142--1144 & WS 5 (FRQ)\\
\textbf{Wed Sep 30} & \textbf{UNIT 3 EXAM (double block)} & all & review sheet & ---\\
Fri Oct 2 & Reteach / float & --- & --- & ---\\
\bottomrule
\end{tabular}

{\sffamily\footnotesize Dates assume continuous instructional weeks; shift
for any co-op breaks.}

\vspace{0.5em}
\section*{\sffamily By the end of this unit you can\ldots}

\begin{itemize}[leftmargin=*,itemsep=0.12em]
  \item Identify every intermolecular force present and rank substances by
        their strength --- including when dispersion beats dipole--dipole. \ced{3.1}
  \item Classify solids (ionic, metallic, covalent network, molecular) and
        predict properties from structure. \ced{3.2}
  \item Explain phase changes and vapor pressure at the particle level. \ced{3.3}
  \item Apply $PV = nRT$ and Dalton's law of partial pressures. \ced{3.4}
  \item State the KMT assumptions and connect them to Maxwell--Boltzmann
        distributions. \ced{3.5}
  \item Predict when and why real gases deviate from ideal behavior. \ced{3.6}
  \item Compute molarity, perform dilutions, and read particulate diagrams
        of solutions. \ced{3.7, 3.8}
  \item Choose a separation technique and justify it from the property it
        exploits. \ced{3.9}
  \item Predict solubility from ``like dissolves like'' with structural
        evidence. \ced{3.10}
  \item Relate wavelength, frequency, and photon energy; match spectroscopy
        type to molecular transition. \ced{3.11, 3.12}
  \item Use $A = \epsilon bc$ and a calibration curve to find an unknown
        concentration. \ced{3.13}
\end{itemize}

\begin{apctrap}[Where students lose points in this unit]
Calling dispersion forces ``weak'' without checking molar mass (\ce{BI3}
outboils \ce{NH3}); forgetting to convert to kelvin in gas problems; saying
real gases deviate ``because they're real'' rather than naming volume and
attraction; treating solubility as a yes/no property instead of a
comparison; and reading a Beer--Lambert calibration curve as though the
slope were the concentration.
\end{apctrap}

\end{document}
